\chapter{Реагування на інциденти та цифрова форензика шкідливого коду}
\chaptershort{Реагування та форензика}
\label{chap:ch16}

Виявлення загрози (розділи~\ref{chap:ch06}--\ref{chap:ch13}) --- не кінець, а початок роботи: за ним має настати \term{реагування на інцидент} --- організований процес стримування, усунення й відновлення, а також \term{цифрова форензика} --- дослідження артефактів для встановлення фактів. Зрілість захисту визначається не лише здатністю виявити, а й здатністю швидко й правильно відреагувати: невміле реагування може завдати більшої шкоди, ніж сам інцидент (знищити докази, не усунути присутність, спричинити повторне зараження). Реагування скорочує час присутності й збитки (розділ~\ref{chap:ch13}), а форензика встановлює повну картину, живить розвідку загроз (розділ~\ref{chap:ch12}) і, за потреби, доказову базу. Розділ розглядає життєвий цикл реагування, його фази, принципи цифрової форензики (докази, їх цілісність, порядок летючості), аналіз інциденту ЗПК, антифорензику й протидію їй, а також нормативні й організаційні аспекти в Україні, і завершується наскрізним кейсом реагування на зараження. Виклад спирається на телеметрію (розділ~\ref{chap:ch10}) й аналіз зразків (розділи~\ref{chap:ch06}--\ref{chap:ch07}).

\section{Життєвий цикл реагування на інцидент}\label{sec:ch16-lifecycle}

Реагування --- не хаотична реакція, а структурований процес. Класична модель~\cite{ch16:nist-ir,ch16:sans-ir} виділяє фази: \emph{підготовка} (процеси, інструменти, навички, телеметрія --- \emph{до} інциденту); \emph{виявлення й аналіз} (підтвердження інциденту, оцінка масштабу); \emph{стримування} (зупинення поширення); \emph{усунення} (видалення присутності загрози); \emph{відновлення} (повернення до нормальної роботи); \emph{уроки} (аналіз після інциденту, вдосконалення). Ці фази утворюють цикл, у якому уроки живлять підготовку до наступних інцидентів.

\begin{tikzfig}
\begin{tikzpicture}[font=\scriptsize,
  b/.style={draw,rounded corners=2pt,minimum height=5.5mm,text width=16mm,align=center,font=\tiny},
  >={Stealth[length=1.4mm]}]
\node[b] (p) at (0,0) {підготовка};
\node[b] (d) at (2.1,0) {виявлення / аналіз};
\node[b] (c) at (4.2,0) {стримування};
\node[b] (e) at (6.3,0) {усунення};
\node[b] (r) at (6.3,-1.3) {відновлення};
\node[b] (l) at (3.15,-1.3) {уроки};
\draw[->] (p)--(d); \draw[->] (d)--(c); \draw[->] (c)--(e); \draw[->] (e)--(r);
\draw[->] (r)--(l); \draw[->] (l) to[out=150,in=-120] (p);
\end{tikzpicture}
\figcap{Рис.~\thechapter.1 --- Життєвий цикл реагування на інцидент як безперервний цикл}
\end{tikzfig}

Критична, часто недооцінена фаза --- \emph{підготовка}: саме вона визначає успіх решти. Без завчасно зібраної телеметрії (розділ~\ref{chap:ch10}), налагоджених процесів, визначених ролей, інструментів і навичок реагування перетворюється на імпровізацію під тиском, коли ціна помилки найвища. Тому зрілі організації інвестують у підготовку й регулярні навчання (кіберполігони, настільні вправи), а не лише в засоби виявлення. Основну мету кожної фази узагальнено в табл.~\thechapter.1.

\tabhead{Таблиця~\thechapter.1 --- Фази реагування та їх основна мета}

\begin{center}\footnotesize
\begin{tabularx}{\linewidth}{L{22mm} Y L{20mm}}
\toprule
Фаза & Основна мета & Ключовий ризик \\
\midrule
Підготовка & готовність процесів і даних & недооцінка \\
Виявлення й аналіз & підтвердження, масштаб & хибна оцінка \\
Стримування & зупинити поширення & втрата доказів \\
Усунення & видалити присутність & неповнота \\
Відновлення & повернути роботу & повторне зараження \\
Уроки & вдосконалення & формальність \\
\bottomrule
\end{tabularx}
\end{center}

\section{Стримування, усунення, відновлення}\label{sec:ch16-containment}

\term{Стримування} --- першочергове зупинення поширення й шкоди: ізоляція уражених вузлів (розділ~\ref{chap:ch10}), блокування зв'язку з командними серверами, відключення скомпрометованих облікових записів. Стримування балансує швидкість і збереження доказів: різке вимкнення зупиняє загрозу, але знищує летючі докази (пам'ять), тому рішення залежить від пріоритетів (негайна зупинка шкоди проти повноти розслідування).

\term{Усунення} видаляє присутність загрози: шкідливі файли, механізми персистентності (розділ~\ref{chap:ch05}), створені облікові записи, бекдори. Ключова вимога --- \emph{повнота}: якщо лишити хоч один механізм персистентності, загроза повернеться. Тому усунення спирається на повне розуміння інциденту (що саме зробив ЗПК), отримане аналізом. \term{Відновлення}~\cite{ch16:recovery} повертає системи до нормальної роботи з контролем, що загрозу усунуто й уразливість, яка дала доступ, закрито (інакше --- повторне зараження). Відновлення з резервних копій вимагає впевненості, що копія не містить тієї самої закладки.

\section{Цифрова форензика: докази та їх цілісність}\label{sec:ch16-forensics}

\term{Цифрова форензика} досліджує цифрові артефакти для встановлення фактів інциденту: що сталося, коли, як, які системи й дані зачеплено, хто (за можливості). Її принципи~\cite{ch16:casey,ch16:carrier} забезпечують, щоб висновки були обґрунтованими, відтворюваними й (за потреби) юридично прийнятними. Міжнародний стандарт ISO/IEC~27037~\cite{ch16:iso27037} визначає настанови з ідентифікації, збору, отримання й збереження цифрових доказів.

Ключові принципи: \emph{цілісність доказів} (докази не повинні змінюватися при зборі --- працюють із копіями, фіксують геші); \emph{ланцюг збереження} (chain of custody --- документований облік, хто, коли й як поводився з доказом); \emph{порядок летючості} (order of volatility~\cite{ch16:orderofvolatility} --- спершу збирають найлетючіші дані). Порядок летючості критичний: регістри й кеш, оперативна пам'ять, мережеві з'єднання зникають за секунди-хвилини, тоді як диск і резервні копії зберігаються довго.

\begin{tikzfig}
\begin{tikzpicture}[font=\scriptsize,
  b/.style={draw,rounded corners=2pt,minimum height=5mm,text width=20mm,align=center,font=\tiny},
  >={Stealth[length=1.4mm]}]
\node[b] (1) at (0,0) {регістри, кеш (секунди)};
\node[b] (2) at (0,-0.95) {оперативна пам'ять, з'єднання};
\node[b] (3) at (0,-1.9) {процеси, тимчасові файли};
\node[b] (4) at (0,-2.85) {диск, журнали};
\node[b] (5) at (0,-3.8) {резервні копії, архіви (довго)};
\draw[->] (1)--(2); \draw[->] (2)--(3); \draw[->] (3)--(4); \draw[->] (4)--(5);
\node[font=\tiny,rotate=90,anchor=south] at (-1.6,-1.9) {черговість збору};
\end{tikzpicture}
\figcap{Рис.~\thechapter.2 --- Порядок летючості: найлетючіші дані збирають першими}
\end{tikzfig}

Для ЗПК особливо цінний аналіз пам'яті (розділ~\ref{chap:ch07}): він розкриває розгорнутий код, ключі, ін'єкції, безфайлові загрози, яких немає на диску. Тому за можливості пам'ять знімають \emph{до} вимкнення вузла --- що входить у суперечність зі швидким стримуванням і вимагає зваженого рішення.

Принцип цілісності доказів заслуговує практичного пояснення. Форензичний аналіз ніколи не ведуть на оригіналі: з носія знімають побітову копію (образ), обчислюють її криптографічний геш і працюють із копією, періодично звіряючи геш, щоб довести незмінність. Будь-яка дія на оригіналі --- навіть просте відкриття файлу --- змінює часові мітки й може знищити докази, тому використовують апаратні чи програмні блокувальники запису. Ланцюг збереження документує кожен перехід доказу з рук у руки з часом, особою й метою, аби в суді чи при службовому розслідуванні можна було довести, що доказ не підмінено й не спотворено. Ці формальності не бюрократія: саме вони відрізняють доказ, на який можна спертися в юридично значущому рішенні (розділ~\ref{chap:ch02}), від припущення, яке опонент легко оскаржить. Порушення цілісності на будь-якому етапі знецінює доказ незворотно --- відновити довіру до нього вже неможливо.

Варто також розрізняти форензику \emph{живої} та \emph{мертвої} системи. Аналіз живої системи (до вимкнення) дає летючі дані --- пам'ять, активні з'єднання, запущені процеси, --- але сам факт роботи аналітика на живій системі вносить зміни, які треба документувати. Аналіз мертвої системи (з образу вимкненого носія) гарантує незмінність, але втрачає летюче. Для ЗПК, особливо безфайлового (розділ~\ref{chap:ch11}), вибір часто на користь живого аналізу з ретельним документуванням, бо втрата пам'яті означає втрату головного доказу. Це рішення --- частина завчасної підготовки, а не імпровізації під час інциденту.

\section{Аналіз інциденту ЗПК}\label{sec:ch16-analysis}

Форензичний аналіз інциденту ЗПК~\cite{ch16:forensic-ua} відновлює повну картину: вектор початкового доступу (як потрапив), механізми персистентності (як закріпився), дії (що робив --- бокове переміщення, збір і вивантаження даних, шкода), індикатори (для виявлення решти уражень і розвідки), масштаб (які системи й дані зачеплено). Аналіз поєднує дослідження зразка (розділи~\ref{chap:ch06}--\ref{chap:ch07}), артефактів системи (журнали, реєстр, файлова система, пам'ять) і телеметрії (розділ~\ref{chap:ch10}) в єдину часову лінію подій (timeline). Основні джерела форензичних артефактів і факти, які вони розкривають, наведено в табл.~\thechapter.2.

\tabhead{Таблиця~\thechapter.2 --- Джерела форензичних артефактів при розслідуванні ЗПК}

\begin{center}\footnotesize
\begin{tabularx}{\linewidth}{L{22mm} Y}
\toprule
Джерело & Що розкриває \\
\midrule
Оперативна пам'ять & розгорнутий код, ключі, ін'єкції, безфайлові загрози \\
Файлова система & зразки, часові мітки, видалені файли \\
Реєстр / автозапуск & механізми персистентності \\
Журнали подій & дії, автентифікації, бокове переміщення \\
Мережеві записи & зв'язок із C2, вивантаження даних \\
\bottomrule
\end{tabularx}
\end{center}

Повноту аналізу зручно характеризувати відновленою часткою ланцюга ураження. Якщо ланцюг має $K$ етапів (розділ~\ref{chap:ch04}), а аналіз достовірно відновив $k$ із них, то повнота розслідування

\begin{equation}
\mathrm{Compl} = \frac{k}{K} \in [0;1],
\label{eq:ch16-completeness}
\end{equation}

причому неповне розслідування (пропущений етап) залишає ризик: невиявлений механізм персистентності чи уражений вузол спричинить повторний інцидент. Формула~\eqref{eq:ch16-completeness} підкреслює, чому мета форензики --- відновити \emph{весь} ланцюг, а не лише очевидний фрагмент: саме прогалини, а не знайдене, визначають залишковий ризик.

\section{Антифорензика та протидія їй}\label{sec:ch16-antiforensics}

Зловмисники застосовують \term{антифорензику}~\cite{ch16:antiforensics} --- приховування й знищення слідів: видалення журналів, затирання файлів, підробку часових міток, безфайлове виконання (немає файлу для аналізу), шифрування. Мета --- ускладнити або унеможливити розслідування, збільшивши час присутності й уникнувши атрибуції.

Протидія антифорензиці спирається на кілька принципів. \emph{Надмірність джерел}: сліди лишаються в багатьох місцях (централізована телеметрія поза ураженим вузлом, мережеві записи, резервні копії), тож знищення одного джерела не знищує всіх. \emph{Завчасний збір}: централізований збір телеметрії (розділ~\ref{chap:ch10}) \emph{до} інциденту зберігає докази поза досяжністю зловмисника на ураженому вузлі. \emph{Аналіз пам'яті}: безфайлові загрози, невидимі на диску, проявляються в пам'яті. \emph{Виявлення самої антифорензики}: ознаки затирання (прогалини в журналах, аномальні часові мітки) самі є індикаторами. Надмірність джерел зручно обґрунтувати кількісно. Якщо один і той самий факт лишає слід у $m$ незалежних джерелах, а зловмисник знищує кожне з імовірністю $q$, то ймовірність, що факт \emph{не} вціліє в жодному,

\begin{equation}
P_{\text{втрати}} = q^{m},
\label{eq:ch16-survival}
\end{equation}

спадає експоненційно зі зростанням кількості незалежних джерел. Формула~\eqref{eq:ch16-survival} формалізує, чому надмірність --- найдієвіша протидія антифорензиці: навіть за високого $q$ (зловмисник майстерно затирає сліди на вузлі) кілька незалежних джерел, зокрема централізована телеметрія поза його досяжністю, роблять повну втрату доказів малоймовірною. Ключовий висновок: централізована, захищена від втручання телеметрія --- найкраща протидія антифорензиці, бо переносить докази туди, куди зловмисник на ураженому вузлі не дістане.

\section{Нормативні й організаційні аспекти в Україні}\label{sec:ch16-ukraine}

Реагування на кіберінциденти в Україні регулюється законодавством (розділ~\ref{chap:ch02}): обов'язок повідомляти про інциденти, взаємодія з CERT-UA (розділ~\ref{chap:ch12}), вимоги до критичної інфраструктури. Стандарти ISO/IEC~27035 (управління інцидентами) та ISO/IEC~27037 (цифрові докази) задають міжнародно узгоджені практики. Організаційно реагування потребує визначеної команди (CSIRT/SOC), ролей, планів, каналів ескалації й взаємодії з правоохоронними органами, коли інцидент має ознаки злочину (ст.~361-1 КК України). Вітчизняні дослідження систематизують організацію реагування відповідно до вимог законодавства~\cite{ch16:ir-ua}. Юридична прийнятність доказів вимагає суворого дотримання цілісності й ланцюга збереження з самого початку, бо порушення на етапі збору знецінює докази незворотно.

\section{Кейс: наскрізне реагування на зараження}\label{sec:ch16-case}

Розглянемо наскрізне реагування на підтверджене зараження (захисна задача, за моделлю NIST~\cite{ch16:nist-ir}).

\emph{Крок 1. Підготовка (до інциденту).} Налагоджено централізовану телеметрію (розділ~\ref{chap:ch10}), визначено команду, ролі, план, інструменти; проведено навчання.

\emph{Крок 2. Виявлення й аналіз.} Сигнал (розділи~\ref{chap:ch10}--\ref{chap:ch11}) підтверджено як інцидент; оцінено попередній масштаб, зафіксовано летючі докази (пам'ять) із дотриманням порядку летючості й цілісності.

\emph{Крок 3. Стримування.} Уражені вузли ізольовано, зв'язок із командними серверами заблоковано, скомпрометовані облікові записи відключено --- із балансом між зупинкою шкоди й збереженням доказів.

\emph{Крок 4. Аналіз і усунення.} Форензичний аналіз відновлює ланцюг ураження~\eqref{eq:ch16-completeness}: вектор доступу, персистентність, дії, масштаб. На основі повного розуміння усуваємо \emph{всі} механізми присутності й закриваємо вразливість доступу.

\emph{Крок 5. Відновлення та уроки.} Відновлюємо системи з перевірених копій, переконуємося у відсутності повторного зараження, видобуваємо індикатори для розвідки (розділ~\ref{chap:ch12}) і нові правила виявлення (розділ~\ref{chap:ch11}), документуємо уроки й удосконалюємо підготовку.

Кейс демонструє, що реагування --- це дисциплінований цикл, у якому підготовка визначає успіх, стримування балансує швидкість і докази, усунення вимагає повноти, а уроки замикають контур удосконалення. Форензика перетворює інцидент на знання, що зміцнює захист проти наступних загроз.

\section*{Висновки до розділу}
\addcontentsline{toc}{section}{Висновки до розділу}

Виявлення загрози --- початок роботи: за ним настає реагування (стримування, усунення, відновлення) і форензика (встановлення фактів). Зрілість захисту визначається здатністю швидко й правильно відреагувати, бо невміле реагування завдає більшої шкоди, ніж сам інцидент. Реагування --- структурований цикл фаз (рис.~\thechapter.1), у якому недооцінена підготовка визначає успіх решти, а уроки живлять підготовку до наступних інцидентів. Стримування балансує швидкість і збереження доказів, усунення вимагає повноти (один лишений механізм персистентності повертає загрозу), відновлення --- впевненості в усуненні й закритті вразливості. Цифрова форензика спирається на цілісність доказів, ланцюг збереження й порядок летючості (рис.~\thechapter.2; ISO/IEC~27037), а для ЗПК особливо цінний аналіз пам'яті. Повнота розслідування~\eqref{eq:ch16-completeness} визначає залишковий ризик: прогалини в ланцюгу ураження спричиняють повторні інциденти. Антифорензику долають надмірністю джерел, завчасним централізованим збором і аналізом пам'яті. В Україні реагування регулюється законодавством і взаємодією з CERT-UA. Головний висновок: реагування й форензика перетворюють виявлений інцидент на кероване усунення й на знання, що зміцнює захист; їх сила --- у підготовці, повноті, цілісності доказів і замкнутому циклі вдосконалення.

\selfcheck
\begin{enumerate}
\item Чому виявлення загрози є початком, а не кінцем роботи? Чим небезпечне невміле реагування?
\item Опишіть фази життєвого циклу реагування (рис.~\thechapter.1). Чому підготовка є критичною й недооціненою?
\item Що таке стримування? Як воно балансує швидкість зупинки шкоди й збереження доказів?
\item Чому усунення вимагає повноти? Що станеться, якщо лишити механізм персистентності?
\item Які умови має задовольняти відновлення, щоб уникнути повторного зараження?
\item Що таке цифрова форензика? Назвіть її ключові принципи (цілісність, ланцюг збереження, порядок летючості).
\item Поясніть порядок летючості (рис.~\thechapter.2). Чому пам'ять знімають до вимкнення вузла?
\item Що відновлює форензичний аналіз інциденту ЗПК? Як поєднуються зразок, артефакти й телеметрія?
\item Поясніть повноту розслідування~\eqref{eq:ch16-completeness}. Чому прогалини визначають залишковий ризик?
\item Що таке антифорензика? Наведіть приклади технік приховування слідів.
\item Якими принципами долають антифорензику? Чому централізована телеметрія є найкращою протидією?
\item Як реагування на кіберінциденти регулюється в Україні? Яка роль CERT-UA та стандартів ISO?
\item Опишіть наскрізне реагування на зараження (кейс підрозділу~\ref{sec:ch16-case}). Як уроки замикають цикл?
\item Дослідницьке: розробіть план реагування й форензичного збору для модельного інциденту ЗПК та оцініть його відповідність ISO/IEC~27037.
\item Дослідницьке: дослідіть ефективність аналізу пам'яті проти набору технік антифорензики на модельних зразках і оцініть залишкові прогалини.
\end{enumerate}
